\section{Selected derivations and assumptions}\label{app:derivations}
\subsection{Population recovery and a calibrated readout}\label{app:enzyme}
The scalar recovery task starts at carrier state 10, targets state 28, and has deadline 120 in the model's units. Capacity lies in $[5/8,9/14]\cup[1,11/8]$ and has mean one. The exact compatible intervals are $[20/41,1]$ from pooled information and $[33/49,4/5]$ after the stated functional readout. Rounded percentages in the main text describe these intervals; calculations use the exact endpoints.

The two 82-unit populations have equal mean because
\[
42\frac9{14}+40\frac{11}{8}=27+55=82.
\]
In the source's separated support, all capacities in the lower band fail the task and all in the upper band complete it. If $p$ is the upper-band mass, its mean can be at most $(1-p)9/14+p\,11/8$. The mean-one constraint therefore gives $p\ge20/41$. The upper bound is one, attained by the uniform population. The kinetic classification, literal source solutions, and attainment of all intermediate fractions are proved in E; the mean calculation alone would not establish task completion.

For the calibrated readout, the largest positive fraction at given $p$ is $p+0.02(1-p)$ and the smallest is $0.9p$. Thus compatibility with $r\in[0.68,0.72]$ requires
\[
0.68\le0.02+0.98p,\qquad0.9p\le0.72,
\]
or $33/49\le p\le4/5$. For any $p$ in this interval the attainable readout interval $[0.9p,0.02+0.98p]$ intersects $[0.68,0.72]$, so suitable $s,f,r$ exist. E supplies a population witness in the same mean-constrained kinetic class. The decision margin is $33/49-2/3=1/147$.

For completeness, the general finite-feature principle behind this example has a short conventional proof \citep{E}. Let population weights $w$ range over a finite simplex, observations be $Aw$, and target be $b^\top w$. Include the normalization row $(1,\ldots,1)$ in $A$. The target is determined by the observations for \emph{every} population precisely when $b^\top$ belongs to the row span of $A$. Sufficiency follows by writing $b^\top=v^\top A$. If the condition fails, linear algebra provides $h\in\ker A$ with $b^\top h\ne0$. The normalization row gives $\sum h_i=0$. For an interior weight vector $w_0$ and sufficiently small $\epsilon>0$, both $w_0\pm\epsilon h$ are populations with the same observations and different targets. This global obstruction does not preclude exact identification at a particular boundary record, or a useful threshold bound without point identification. This argument is conventional; we do not represent it as a newly compiled Lean theorem.

\subsection{Dilution, physical occupancy, and the error budget}\label{app:hook}
Write the capture and detector occupation fractions as $p_C,p_D$. The independent-site equilibrium source satisfies
\[
C=up_C+\frac{Kp_C}{1-p_C},\quad
D=up_D+\frac{Jp_D}{1-p_D},\quad B(u)=up_Cp_D.
\]
For $C,K>0$, the physical root is
\[
p_C=\frac{2C}{u+C+K+\sqrt{(u-C)^2+2K(u+C)+K^2}}.
\]
The same formula with $D,J$ gives $p_D$. This is the formula used in Figure~\ref{fig:hook}; it satisfies the reagent balances and gives $B(u)\sim CD/u$ at high concentration. At unit parameters, $B(0.08)=B(50)$ by the source involution. The plotted equality is illustrative of the exact structure, not a numerical proof of it.

With observations $Y_1=gB(u)+e_1$ and $Y_d=grB(u/d)+e_d$, the contrast is $Z=g\Delta+e_d-e_1$. Let $|e_1|\le\varepsilon_1$, $|e_d|\le\varepsilon_d$, and $g\ge g_{\min}>0$. Source bounds $\Delta\le0$ on one class and $\Delta\ge m$ on the other imply
\[
Z\le\varepsilon_1+\varepsilon_d\quad\hbox{on the low class},\qquad
Z\ge g_{\min}m-\varepsilon_1-\varepsilon_d\quad\hbox{on the high class}.
\]
The two ranges are strictly separated when $2(\varepsilon_1+\varepsilon_d)<g_{\min}m$. Taking $g_{\min}=1$, $m=1/25$, and equal errors gives $\varepsilon<1/100$. Source I proves the continuum margin over $C,D,K,J\in[9/10,11/10]$, $d\in[9,11]$, and $r\in[19/20,21/20]$ by exact polynomial certificates. A plotted unit-parameter curve is not a substitute for those certificates. The availability model treats the inaccessible fraction as invisible to both sites; partially masked material that consumes one reagent requires a different source model.

\subsection{A joint reporting event, not a posterior}\label{app:specimen}
Let $N$ be the original-specimen target count, $a,b$ the unique effective fractions examined, and $X,Y$ the two state-dependent recovery probabilities, with $a+b\le1$. Conditional independence of target paths gives negative probability $\E[(1-aX-bY)^N]$. The preparation-state average is taken after the power because the state is shared; replacing $X,Y$ with their means generally does not preserve this probability.

For balanced unique effective fractions $a=b$, count threshold $K$, and $m\ge1$ independent reference pairs, set $H=(1-a)^K$. Study S reduces the sharp false-issuance bound to
\[
G_m(H)=\max_{0\le z\le1}z^m\{1-(1-H)z\}.
\]
The reduction uses the full recovery-law optimization and the paired-reference experiment; it is not obtained merely by guessing a scalar recovery probability. Differentiating the final polynomial gives
\[
G_m(H)=
\begin{cases}
H,&H\ge1/(m+1),\\[2pt]
\displaystyle\frac{m^m}{(m+1)^{m+1}(1-H)^m},&H<1/(m+1).
\end{cases}
\]
At $a=9/20$, $K=6$, $m=9$, $H=(11/20)^6=0.027680640625$ and $G_9(H)\approx0.0498772$. The exact rational value is reproduced in the calculation file. The limiting floor $H$ has a material origin: even favorable recovery cannot find a target in an unexamined part of a specimen.

Let $G$ denote gate passage and $N_-$ the negative record. The theorem bounds $\Prb(G\cap N_-)$ for each admitted law when the target count is at least $K$. In contrast, $\Prb(N_-\mid G)=\Prb(G\cap N_-)/\Prb(G)$, when the denominator is positive. A joint bound alone does not give the same conditional bound. A posterior probability of count below $K$ would additionally require a prior and a likelihood model. These distinctions do not diminish the usefulness of the reporting rule; they specify the experiment for which its reliability is proved.

\subsection{Complete paths across the full coverage range}\label{app:paths}
The uniform type-level conditions in V are $0\le x,y,a,b\le1$, $x+a\ge c$, $y+b\ge c$, and $|x-y|\le d$, with $0\le c\le2$ and $d\ge0$. Here $x,y$ are the successive stage probabilities in the first condition, and $a,b$ those in the second. The exact balanced floor is
\[
R(c,d)=\frac{c^2-t^2}{4},\qquad t=\min\{d,c,2-c\}.
\]
To see the geometry, first reduce the cross-condition sums to $c$ by clipping the first-condition probabilities into $[\max(0,c-1),\min(c,1)]$. This can only reduce the complete-path lower expression and cannot increase the distance between the clipped values $X,Y$. The remaining algebra is
\[
\frac{XY+(c-X)(c-Y)}2
=\frac{c^2-(X-Y)^2+(X+Y-c)^2}{4}\ge R(c,d).
\]
Taking $X=(c+t)/2$, $Y=(c-t)/2$, $a=Y$, $b=X$ attains the bound. The choices lie in $[0,1]$ because $t\le c$ and $t\le2-c$. In particular, for $c>1$, $R(c,d)\ge c-1$.

With recorded mean path probabilities $u,v$ satisfying $(u+v)/2\ge g$, recoverable fraction $p\ge\theta$, and $n=2m$ independent units split equally, the all-negative probability is $(1-pu)^m(1-pv)^m\le(1-\theta g)^n$. If this holds in every valid calibration environment and invalid environments have probability at most $\delta$, the full bound is $\delta+(1-\delta)(1-\theta g)^n$. Validity conditional on the calibration environment is essential. V also proves a mean-square mismatch extension; it should not be confused with a bound inferred from pooled stage averages alone.

\subsection{The two-founder calculation and the loading floor}\label{app:counts}
At $T=\log2$, a fast Yule family has $\Prb(Z=j)=2^{-j}$ for $j\ge1$, while a slow family stays at one. Each founder's marginal class is half slow and half fast. Independent classes give mixture weights $1/4,1/2,1/4$ for two slow, mixed, and two fast families; the shared class gives weights $1/2,0,1/2$. Hence, for $k\ge2$,
\[
F_I(k)=1-\frac{k+5}{4\,2^k},\qquad
F_W(k)=1-\frac{k+1}{2^{k+1}}.
\]
At $k=6$, these are $245/256$ and $121/128$; at $k=7$, $125/128$ and $31/32$. A single chronological family record has the same mixture law in both preparations. The difference enters only when multiple founders are combined. C's formal development verifies the relevant chronological source calculation and finite cutoffs; the general formula is also derived conventionally in that manuscript.

For the separate amplification model A, let $M$ be a Poisson-loaded initial target count with mean $\lambda$. Conditional on $M=0$, its observation law is the blank law. Any classifier with blank false-positive probability at most $\alpha$ therefore has miss probability at least $\Prb(M=0)(1-\alpha)=e^{-\lambda}(1-\alpha)$. This loading floor is distinct from the stronger full-timing obstruction in the five-state example. It applies to the stated reaction-level loading model, not automatically to the original specimen.

\subsection{The sharp two-pool material account}\label{app:material}
Let $q_1,q_2$ be lower bounds on the two collections, $B$ an upper bound on initial mobile-plus-reserve inventory, $J$ an upper bound on reserve immediately before recovery, and $H$ an upper bound on material introduced during recovery. The retained mobile and reserve fractions are bounded by $e,s$, with $0\le e\le s\le1$. The exact minimum fresh entry over the admitted material histories is
\begin{equation}\label{eq:material}
\begin{split}
L=\max\{&0,\ q_1-B,\ q_1+q_2-B-H,\\
&e q_1+q_2-eB-(s-e)J-H,\ s q_1+q_2-sB-H\}.
\end{split}
\end{equation}
Let $f_1,f_2$ be fresh entries in the two windows and $P,R$ the mobile and reserve amounts just before recovery. The source balances imply
\[
q_1+P+R\le B+f_1,\qquad
q_2\le eP+sR+H+f_2,\qquad R\le J,
\]
with all material quantities nonnegative and $0\le e\le s\le1$. Each term in \eqref{eq:material} is a lower bound on $f_1+f_2$. For example, multiply the first inequality by $e$, add the second, and use $(s-e)R\le(s-e)J$ and $ef_1\le f_1$ to obtain the reserve-sensitive term. Replacing $eP+sR$ by $s(P+R)$ gives the final term. Bounding retention by one gives the total-input term.

Define $A=\pos{q_1-B}$ and $x=\pos{B-q_1}$. The largest retained old stock compatible with the leftover amount is
\[
K_{\rm old}=ex+(s-e)\min(J,x),
\qquad L=A+\pos{q_2-H-K_{\rm old}}.
\]
If $q_1\le B$, collect $q_1$ from old material and place $\min(J,x)$ of the remaining $x$ in reserve, the rest in the mobile pool. If $q_1>B$, use $q_1-B$ fresh units to complete that collection and leave no old material. In either case, the second window uses retained old stock, at most $H$ recovery input, and exactly the positive remaining deficit as fresh entry; unneeded material can be lost. The source permits the required transfers and selective collection. This constructs an attaining account and explains the equivalence with the maximum of five expressions. It does not impose a species-specific finite rate on those transfers.

For the washed example, exact rational substitution yields the five terms $0,-4.2,-0.6,1.69,-0.18$. For the unwashed example, use $e=s=0.9$ and $q_2=5.5$, giving $L=1.52$. A common witness has $q_1^{\rm true}=6$, pre-recovery mobile stock 2 and reserve 2, recovery input 0.2, and second-window fresh entry 2.1. Its second collections are $0.05(2)+0.9(2)+0.2+2.1=4.2$ and $0.9(4)+0.2+2.1=5.9$, each at the upper edge of the corresponding error interval. Both records are thus physically matched within the admitted source.

In the washed record family, the active reserve-sensitive expression is $3.39-0.85J$. Requiring it to reach 1.6 gives $J\le179/85\approx2.1059$. This calculation retains the other premises of the worked family. It must not be interpreted as a universal reserve threshold.

An initial reserve bound $\overline R$ and gross uptake bound $\overline T$ imply only $R\le\overline R+\overline T+f_1$, not necessarily $R\le\overline R+\overline T$. Nevertheless substitution $J_0=\overline R+\overline T$ in the same lower-bound formula is sound. Indeed the reserve-sensitive derivation then has coefficient $s$ on $f_1$, which is at most one. The other four inequalities remain valid. The companion proves attainment in this uptake-bounded class as well. This distinction avoids treating fresh reserve material as if it could not exist.

\subsection{Reporter loss, storage, and convergence}\label{app:reporter}
The complete-recovery cost identity in the reporter-loading model is
\begin{equation}\label{eq:cost}
\frac{D_\infty}{q_\infty}=\frac{k}{p+k}\left(1+\frac{p}{c}\right),\qquad q_\infty>0,
\end{equation}
where $p$ is regeneration, $k$ target use, and $c$ reporter catalytic release. Positive coefficients, matched preparation, and finite total association exposure are assumed. The source also constructs identical reporter observations with stationary target fluxes $C_0/6$ and $C_0/2$, where $C_0$ is total cofactor, by redistributing consumption between the target and an unobserved background channel.

Let $x_0$ be reduced cofactor on the matched reporter-free trajectory, $x$ its reported counterpart, $b$ bound cofactor, and $e=x_0-x$. With regeneration coefficient $p$, target coefficient $k$, and catalytic release $c$, source subtraction gives
\[
e'=-(p+k)e+b'+(p+c)b.
\]
For matched preparation $e(0)=b(0)=0$, integration yields
\[
(p+k)D(T)=k\{(p+c)\int_0^T b(t)\,dt+b(T)-e(T)\},
\quad D(T)=k\int_0^T e(t)\,dt.
\]
The conventional existence, comparison, and convergence arguments in R justify passing to complete recovery for its finite-exposure association schedules. Since reporter product satisfies $q_\infty=c\int_0^\infty b(t)\,dt$, the boundary terms vanish and \eqref{eq:cost} follows. The finite integral account and cost algebra are Lean-checked; the statement about every admitted ODE trajectory also depends on those conventional analytic arguments. Reporting only the algebraic declaration as a fully formalized ODE theorem would overstate the evidence.

\section{Source and proof scope}\label{app:map}
The following map names the main local proof modules supporting the selected companion claims. Lean is a proof assistant that checks a stated mathematical implication against its assumptions \citep{lean}; it does not establish the biological truth of those assumptions. Paths are relative to the repository's \texttt{proofs/} directory. The letters below identify enzyme recovery (E), immunoassay (I), specimen recovery (S), recovery paths (V), small populations (C), amplification (A), microbial conversion (M), and reporter loading (R). They are used here for provenance rather than as biological nomenclature. The map concerns specific implications, not the formalization of every manuscript. The source package records source hashes and audits matching strict receipts.

\begingroup\small
\begin{longtable}{@{}>{\raggedright\arraybackslash}p{0.07\linewidth}>{\raggedright\arraybackslash}p{0.43\linewidth}>{\raggedright\arraybackslash}p{0.44\linewidth}@{}}
\toprule Source & Selected module and declaration & Evidence boundary\\\midrule\endhead
E & \path{G6PDReserve/PopulationRoot.lean}: \path{population_root}, \path{sharp_bulk}, \path{sharp_repaired}. & Literal scalar solutions, pooled equivalence, and exact population ranges. The general finite-feature criterion is conventional.\\[4pt]
I & \path{SandwichImmunoassay/Publication.lean}: \path{publication_main}; \path{GeneralBox.lean}; \path{Obstructions.lean}. & Source certificates and bounded decisions; general shape and incubation analysis also use conventional proofs.\\[4pt]
S & \path{SpecimenReliability/AtLeastOneGate.lean}; \path{GateTradeoff.lean}; \path{PolicyComparison.lean}: \path{worked_bounds}. & Source-law reduction and joint-event bounds, including the worked design. Biological transport and reference matching remain assumptions.\\[4pt]
V & \path{FunctionalViability/ExtendedCoverage.lean}: \path{extended_minimax}, \path{moment_calibrated_certificate}. & Full-range path bounds and finite source composition. Type representativeness and conditional calibration must be supplied.\\[4pt]
C & \path{MemoryPrediction/YuleAmbiguity.lean}: \path{chronological_single_observation_equal}, \path{chronological_minimal_endpoints}. & Implemented chronological source and finite cutoffs; the general-$k$ generating-function calculation is conventional.\\[4pt]
A & \path{AssayInformation/Main.lean}: \path{AssayInformation.main}. & All-timing-classifier obstruction and constructed identity observation. Optimal-deadline analysis and physical reporter realization are separate.\\[4pt]
M & \path{MicrobialFunctionAssay/Sharpness.lean}: \path{lower_is_minimum}, \path{uptake_source_bound}, \path{uptake_bound_attained}. & Sharpness in the stated material-history classes. Actual microbial kinetics and inventory calibration are not supplied.\\[4pt]
R & \path{BiochemicalReadout.lean}: \path{finite_account}, \path{turnover_cost}, \path{joint_protocol}; \path{BiochemicalReadoutExtensions.lean}. & Integral algebra, decision constants, and selected recovery inequalities. ODE existence, comparison, convergence, and biochemical mapping have separate roles.\\\bottomrule
\end{longtable}
\endgroup

The reproducible calculations verify arithmetic, material witnesses, and the plotted equilibrium source; they do not re-prove all companion results. Unpublished companion references are identified by their full titles and local version provenance, without invented journal details, author names, or public identifiers. The literature cited for assay practice supplies context and methods, not empirical validation of the synthetic numerical examples.
